\section{Coefficient fields of exact sum-of-squares certificates}
\label{sec:fields}

Let $(f,A)$ be a certified quartic in the sense of
Section~\ref{sec:heights-conventions} whose minimum value is zero, attained
at $p$. Convexity of $f$ is then certified by short rational data, and
Lemma~\ref{lem:taylor-sos} turns the Hessian certificate into a sum of squares
for $f$ itself. The squares produced in this way have coefficients in the
field $\Q(p)$ generated by the coordinates of $p$, which can have large
degree (Section~\ref{sec:algebraic}). This section asks how much of that
field an exact certificate of the optimal value must use. Three questions
are natural.
\begin{enumerate}[label=(\roman*)]
\item Does a rational convexity certificate guarantee a sum of squares of
  \emph{rational} polynomials for $f$? This is the question of rational
  descent.
\item If it does not, which real fields carry an exact sum-of-squares or
  positive semidefinite Gram certificate of $f$, and can the smallest such
  field be forced to be large while the input stays short?
\item Which compact exact certificates remain available when every
  polynomial certificate needs a large field?
\end{enumerate}
The answers matter for exact certification. Exact optimality certificates
are typically obtained by rounding a numerical semidefinite solution to
rational data and verifying a polynomial identity
\cite{PeyrlParrilo2008,KaltofenLiYangZhi2012}. When the optimal value is
attained at an irrational point, such rounding must fail if no rational
certificate exists, and an algebraic certificate is only useful if its
field is of manageable size. The results below separate the certification
of convexity from the certification of the optimal value, and they identify
which certificate representations are long and which stay short.

The main results are the following.
\begin{enumerate}[label=(\alph*)]
\item For every prime $\ell\ge5$ there is an integer quartic in
  $(\ell+1)/2$ variables with an integer positive definite full Hessian Gram;
  each of its coefficients and each Gram entry has $O(\log\ell)$ bits. A sum
  of squares or a positive semidefinite Gram of this quartic over a real
  field $E$ exists exactly when $2^{1/\ell}\in E$
  (Theorem~\ref{thm:fields-prime}). For $\ell=5$ this is a
  ternary quartic, and three variables are the minimum for a failure of
  rational descent among certified quartics with minimum zero
  (Corollary~\ref{cor:fields-dimension}); an explicit ternary instance has 31
  monomials and coefficients of absolute value at most $448$
  (Example~\ref{ex:fields-ternary}).
\item For every $k$ there is a rational quartic in $3k$ variables, with input
  and full Hessian Gram of size polynomial in $k$, whose least real
  certificate field is $\Q(2^{1/5^k})$, of degree $5^k$
  (Theorem~\ref{thm:fields-tower}). Every algebraic certificate even has one
  individual coefficient of degree at least $5^k$
  (Theorem~\ref{thm:fields-individual}). Shared root circuits and rational
  certificates with auxiliary root variables are nevertheless of polynomial
  size (Theorem~\ref{thm:fields-compact}).
\item At the tower point, the rational quadratic relations, the span of their
  products, and the rational quartics with zero value and gradient can be
  determined explicitly. These descriptions yield an exact rational recognition
  algorithm (Theorem~\ref{thm:fields-tower-spaces} and
  Corollary~\ref{cor:fields-recognition}).
\item Conversely, if the rational quadratics vanishing at $p$ form a space of
  the minimum possible dimension $n+1$ and their products exhaust the
  stationary quartics, then every globally convex rational quartic with
  minimum zero at $p$ and positive definite Hessian there is a rational sum
  of squares, provided some globally convex rational sum of squares of
  degree four vanishes at $p$ with positive definite Hessian there
  (Theorem~\ref{thm:fields-descent}).
\end{enumerate}
All statements concern polynomial sums of squares and polynomial Gram
matrices of the attained value. Rational-function certificates, multipliers,
and block-structured certificates are treated in
Section~\ref{sec:fields-certificates}. No result of this section is a lower
bound for the complexity of a decision problem or for approximation.

\subsection{Certificate models and basic tools}
\label{sec:fields-models}

We use the conventions of Section~\ref{sec:heights-conventions}: $z(X)$ is the
vector of monomials of degree at most two with the constant monomial first,
and $w(X,v)=(v,X\otimes v)$ is the full Hessian basis. For a polynomial $P$
and a monomial $m$, $[m]P$ denotes the coefficient of $m$ in $P$. For a
subfield $E\subseteq\R$, a point $p\in\R^n$, and $d\ge1$ put
\[
 \mathcal I_d^E(p)=\{g\in E[X]_{\le d}:g(p)=0\},\qquad
 \mathcal J_d^E(p)=\{g\in E[X]_{\le d}:g(p)=0,\ \nabla g(p)=0\},
\]
and let $\mathcal W^E(p)$ be the $E$-span of all products $gh$ with
$g,h\in\mathcal I_2^E(p)$. The product rule gives
$\mathcal W^E(p)\subseteq\mathcal J_4^E(p)$.

\begin{definition}[Certificates over a real field]
\label{def:fields-certificates}
Let $E\subseteq\R$ be a subfield and let $f\in\R[X]$ have degree at most
four.
\begin{enumerate}[label=(\alph*)]
\item $f$ is a \emph{sum of squares over $E$} if $f=\sum_jg_j^2$ for finitely
  many $g_j\in E[X]$.
\item $f$ has a \emph{positive semidefinite Gram matrix over $E$} if
  $f=z^{\mathsf T}Qz$ for a symmetric matrix $Q$ with entries in $E$ that is
  positive semidefinite as a real matrix.
\item A subfield $E_0\subseteq\R$ is the \emph{least real field} for one of
  these certificate types if, for every subfield $E\subseteq\R$, a certificate
  of that type over $E$ exists exactly when $E_0\subseteq E$.
\end{enumerate}
\end{definition}

The field $E$ is arbitrary: it need not be a number field and may contain
transcendental elements. A sum of squares over $E$ gives a positive
semidefinite Gram matrix over $E$, namely the sum of the outer products of
the coefficient vectors of the square factors. Over $\Q$ the converse holds,
by rational elimination and the splitting of positive rationals into rational
squares (Lemma~\ref{lem:models-rational-squares}). Over a general real field
the converse fails: Remark~\ref{rem:heights-gram-field} gives a positive
semidefinite Gram over $\Q(\sqrt2)$ of a polynomial that is not a sum of
squares over $\Q(\sqrt2)$; see also \cite[Lemma~1.6 and
Remark~1.7]{ChuaPlaumannSinnVinzant2017}. Every least-field statement below
therefore proves necessity separately for square factorizations and for
positive semidefinite Gram matrices, and proves sufficiency by exhibiting
actual squares.

\begin{lemma}[Basic tools]
\label{lem:fields-tools}
Let $E\subseteq\R$ be a subfield and $p\in\R^n$.
\begin{enumerate}[label=(\alph*)]
\item \emph{Degree.} If $f=\sum_jg_j^2$ with real $g_j$ and $\deg f\le4$,
  then every $g_j$ has degree at most two.
\item \emph{Kernel restriction.} Suppose $\deg f\le4$ and $f(p)=0$. Every
  factor of a real sum of squares of $f$ vanishes at $p$. If
  $f=z^{\mathsf T}Qz$ with $Q\succeq0$, then $Qz(p)=0$. If moreover $Q$ has
  entries in $E$ and $g=(g_1,\ldots,g_t)^{\mathsf T}$ lists a basis of
  $\mathcal I_2^E(p)$, then $f=g^{\mathsf T}Sg$ for a positive semidefinite $S$
  with entries in $E$.
\item \emph{Obstruction functionals.} Let $\varphi$ be a real linear
  functional on $\R[X]_{\le4}$ with $\varphi(g^2)\ge0$ for every
  $g\in\mathcal I_2^E(p)$. If $\deg f\le 4$, $f(p)=0$, and $\varphi(f)<0$,
  then $f$ is neither a sum of squares over $E$ nor has a positive
  semidefinite Gram matrix over $E$.
\item \emph{Odd radicals over real fields.} Let $D\ge1$ be odd and
  $a=2^{1/D}>0$. The minimal polynomial of $a$ over $E$ is $T^r-a^r$ for some
  divisor $r$ of $D$. In particular, if $D=\ell$ is prime and $a\notin E$,
  then $T^\ell-2$ is irreducible over $E$; and if $D=5^k$ with $k\ge1$ and
  $a\notin E$, then $[E(a):E(a^5)]=5$.
\end{enumerate}
\end{lemma}

The proof is in Appendix~\ref{app:fields-tools}. Part (d) is the reason why
the least fields below can be described for \emph{all} real fields, including
transcendental ones: the only real root of unity of odd order is one, so the
constant term of a factor of $T^D-2$ over a real field is a power of $a$.
Sufficiency always comes from Lemma~\ref{lem:taylor-sos}(e): if $f$ has a
rational positive semidefinite full Hessian Gram and a zero $p$ with
$\nabla f(p)=0$, then $f$ is a sum of squares of quadratics whose
coefficients are rational polynomials of degree at most two in the
coordinates of $p$, hence a sum of squares over every subfield $E\subseteq\R$
with $p\in E^n$. The scalar identity
$\int_0^1(1-t)(U+tV)^2dt=2\bigl((U+V/3)/2\bigr)^2+(V/6)^2$ behind it produces
actual squares and introduces no square roots. For a rational $A\succ0$ of
order $h$ we write $\mu_A=\det A/(\tr A)^{h-1}$; then $A\succeq\mu_AI$
(Lemma~\ref{lem:models-det-trace}).

\subsection{Rational convexity certificates do not give rational descent}
\label{sec:fields-prime}

Let $\ell\ge5$ be prime, put $s=(\ell-1)/2$ and $n=s+1$, and use variables
$x_0,\ldots,x_s$. Let $a=2^{1/\ell}$ and
\[
 p=(a^s,a^{s+1},\ldots,a^{2s})\in[1,2)^n .
\]
Its coordinates generate $\Q(a)$, since $a=2/p_s$, a field of degree $\ell$
by Eisenstein's criterion. Define the rational linear functional
\begin{equation}
 \varphi_\ell(P)=[x_s]P+[x_0^2]P .
 \label{eq:fields-prime-functional}
\end{equation}

\begin{lemma}[A coefficient obstruction at $p$]
\label{lem:fields-prime-functional}
Let $E\subseteq\R$ be a subfield with $a\notin E$, and let
$g=g_0+\sum_iv_ix_i+\sum_{i\le j}h_{ij}x_ix_j\in E[X]_{\le2}$ vanish at
$p$. Then $v_s+h_{00}=0$ and $\varphi_\ell(g^2)=v_0^2$.
\end{lemma}

\begin{proof}
A monomial of degree at most two evaluates at $p$ to $1$, to $a^{s+i}$ with
$0\le i\le s$, or to $a^{2s+i+j}$ with $0\le i\le j\le s$. Reduce the
exponents with $a^\ell=a^{2s+1}=2$. The exponents $2s+i+j$ with $i+j\ge1$ lie
in $[2s+1,4s]$ and reduce to $2a^{i+j-1}$ with $i+j-1\le2s-1$. Hence the only
monomials contributing to $a^{2s}$ are $x_s$ and $x_0^2$, and the coefficient
of $a^{2s}$ in $g(p)$ is $v_s+h_{00}$. By Lemma~\ref{lem:fields-tools}(d) the
powers $1,a,\ldots,a^{2s}$ are linearly independent over $E$, so $g(p)=0$
forces $v_s+h_{00}=0$. In $g^2$ the monomial $x_s$ has coefficient
$2g_0v_s$ and $x_0^2$ has coefficient $v_0^2+2g_0h_{00}$. Therefore
$\varphi_\ell(g^2)=v_0^2+2g_0(v_s+h_{00})=v_0^2$.
\end{proof}

The lemma makes $\varphi_\ell$ nonnegative on squares of $E$-quadratics
vanishing at $p$, but only those: $\varphi_\ell((1-x_s)^2)=-2$. The
construction therefore builds a strictly convex rational sum of squares from
vanishing quadratics with no linear $x_0$ term, and subtracts the square of
the vanishing quadratic $r=x_1x_s-2x_0$, which has $\varphi_\ell(r^2)=4$.
Appendix~\ref{app:fields-prime} gives the construction and its quantitative
convexity proof.

\begin{theorem}[Prime least fields]
\label{thm:fields-prime}
For every prime $\ell\ge5$ one can construct, in time polynomial in $\ell$,
an integer quartic $f_\ell$ in the $n=(\ell+1)/2$ variables $x_0,\ldots,x_s$
and an integer positive definite full Hessian Gram of $f_\ell$ with the
following properties.
\begin{enumerate}[label=(\alph*)]
\item $f_\ell$ has $O(\ell^2)$ monomials, and its coefficients and the
  entries of its Hessian Gram have $O(\log\ell)$ bits.
\item $\nabla^2f_\ell\succeq I$ on $\R^n$, $f_\ell(p)=0$, and
  $f_\ell(x)\ge\frac12\norm{x-p}^2$. Thus $p$ is the unique minimizer and
  $\min f_\ell=0$.
\item $\varphi_\ell(f_\ell)=-4$.
\item For every subfield $E\subseteq\R$ the following are equivalent:
  $2^{1/\ell}\in E$; $f_\ell$ is a sum of squares over $E$; $f_\ell$ has a
  positive semidefinite Gram matrix over $E$.
\end{enumerate}
\end{theorem}

\begin{proof}[Proof of (d) from (a)--(c)]
If $a\in E$ then $p\in E^n$, and Lemma~\ref{lem:taylor-sos}(e) applies,
because $f_\ell$ has a rational positive definite Hessian Gram and a
stationary zero at $p$; a sum of squares over $E$ gives a positive
semidefinite Gram over $E$. If $a\notin E$, Lemma~\ref{lem:fields-prime-functional}
shows $\varphi_\ell(g^2)=v_0^2\ge0$ for every $g\in\mathcal I_2^E(p)$, and
Lemma~\ref{lem:fields-tools}(c), together with $\varphi_\ell(f_\ell)=-4$,
excludes both certificate types.
\end{proof}

For the Gram case the argument can be made entirely explicit: if
$f_\ell=z^{\mathsf T}Qz$ with $Q\succeq0$ over $E$, then the constant row of
$Q$ is a vanishing quadratic, so $Q_{1,x_s}+Q_{1,x_0^2}=0$, and
\[
 \varphi_\ell(f_\ell)=2Q_{1,x_s}+Q_{x_0,x_0}+2Q_{1,x_0^2}=Q_{x_0,x_0}\ge0 .
\]
No factorization of $Q$ over $E$ is used.

Theorem~\ref{thm:fields-prime} says more than that rational descent fails.
The least real field for either certificate is exactly $\Q(2^{1/\ell})$;
every real number field carrying a certificate contains it and so has
degree divisible by $\ell$. In particular no finite tower of quadratic
extensions suffices. The field degree grows only linearly with the number of
variables here, and the field has a short radical description, so this is
not an output-size lower bound; Section~\ref{sec:fields-tower} provides one.

\begin{corollary}[Rational lower bounds do not attain the minimum]
\label{cor:fields-nonattainment}
Let $(f,A)$ be a certified quartic with minimum zero that is not a sum of
squares of rational polynomials, for instance $f_\ell$. Then $f+\eta$ has a
rational positive definite Gram matrix for every rational $\eta>0$, and
\[
 \sup\{\gamma\in\Q:f-\gamma\text{ is a sum of squares in }\Q[X]\}=0
\]
is not attained.
\end{corollary}

\begin{proof}
Let $\mu=\mu_A$ and let $p$ be the minimizer. The quantity
$f(a)+\eta-\norm{\nabla f(a)}^2/(2\mu)$ tends to $\eta$ as $a\to p$, so it
is positive for some rational $a$. Lemma~\ref{lem:taylor-sos}(c), applied to
$f+\eta$ at this center, gives a rational positive definite Gram matrix, and
Lemma~\ref{lem:models-rational-squares} gives rational squares. Every
negative rational $\gamma$ is therefore feasible, $\gamma=0$ is excluded by
hypothesis, and positive $\gamma$ are excluded by evaluating at $p$.
\end{proof}

The real semidefinite program is not at fault: it attains the value zero,
with an algebraic optimal Gram (Theorem~\ref{thm:heights-moment}). The
failure is arithmetic. A procedure that rounds near-optimal Gram matrices to
rational data obtains certified lower bounds $-\eta$ for every $\eta>0$, and
never the value zero. Corollary~\ref{cor:fields-nonattainment} makes no
claim about the bit length of these certificates as $\eta\to0$.

The construction behind Theorem~\ref{thm:fields-prime} uses uniform constants
and therefore large coefficients. A hand-sized example with the same
mechanism is the following.

\begin{example}[An explicit ternary quartic]
\label{ex:fields-ternary}
In the variables $(x,y,z)$ put
\[
 r_0=2-2xy,\quad r_1=2x^2-2yz,\quad r_2=2y^2-4z,\quad r_3=2z^2-x,\quad
 r_4=2xz-y,
\]
$A_0=4r_0+5r_1+3r_2+9r_3$, and
\begin{equation}
 F_\star=A_0^2+r_0^2+r_1^2+r_2^2+r_3^2-r_4^2 .
 \label{eq:fields-ternary}
\end{equation}
This integer quartic has 31 monomials and coefficients of absolute value at
most $448$. All five $r_i$ vanish at $p_\star=(a^4/2,a,a^2/2)$, where
$a=2^{1/5}$. Appendix~\ref{app:fields-ternary} prints a rational full
Hessian Gram $M_\star$ of $F_\star$, specified exactly through rational
quadratic data at the center $(3/4,1,1/2)$, and the twelve positive leading
principal minors of the integer matrix $2(M_\star-I)$. They prove
$\nabla^2F_\star\succeq I$ by Sylvester's criterion. Hence $p_\star$ is the
unique minimizer and $\min F_\star=0$. With
$\varphi_\star(P)=2[z]P+4[y^2]P$ one has
\[
 \varphi_\star(r_ir_j)=4\quad\text{if }i=j=4,\qquad
 \varphi_\star(r_ir_j)=0\quad\text{otherwise},
\]
so $\varphi_\star(F_\star)=-4$. For every real field $E$ with $a\notin E$
the $r_i$ form a basis of $\mathcal I_2^E(p_\star)$, so
$\varphi_\star(g^2)\ge0$ on that space, and Lemma~\ref{lem:fields-tools}
gives the same characterization as Theorem~\ref{thm:fields-prime}(d): a sum
of squares or a positive semidefinite Gram over $E$ exists exactly when
$2^{1/5}\in E$. The fifteen products $r_ir_j$ are linearly independent, so
$F_\star$ has the unique Gram matrix $(I_4+gg^{\mathsf T})\oplus(-1)$,
$g=(4,5,3,9)^{\mathsf T}$, on the basis $(r_0,\ldots,r_4)$. Only the
convexity certificate is a finite computation. Given it, the field statements
are proved in Appendix~\ref{app:fields-ternary}; sufficiency uses the rational
Hessian Gram that the certificate provides.
\end{example}

The two examples differ from earlier failures of rational descent in the
combination of hypotheses, not in the failure itself. Scheiderer proved
that rational forms that are sums of squares over $\R$ need not be sums of
squares over $\Q$, and that every prescribed real number field can be
excluded as a field of coefficients by a rational nonnegative ternary quartic
form \cite[Theorem~2.1 and Corollary~2.11]{Scheiderer2016}. Descent does
hold from totally real number fields \cite[Theorem~1.4]{Hillar2009},
\cite[Theorem~1.2]{Scheiderer2016}, and for rational polynomials with a
positive definite real Gram matrix \cite[Theorem~1.2]{Hillar2009}; exact
rounding methods work in that interior case
\cite{PeyrlParrilo2008,DavisPapp2022}. A quaternary
quartic form that is a sum of squares over $\Q(2^{1/3})$ but not over $\Q$ is
attributed to Capco, Laplagne, and Scheiderer and proved in
\cite[Section~3.1]{Laplagne2024}; that form is not convex. What
Theorem~\ref{thm:fields-prime} adds is the simultaneous presence of a strict
rational full Hessian certificate, global strong convexity, an attained
rational minimum, small integer data, and an exact description of all real
fields carrying either certificate type, with an unbounded prime degree. The
abstract mechanisms, kernel restriction of Gram matrices at real zeros and
odd-degree field arithmetic, are classical.

\begin{corollary}[Three variables are needed]
\label{cor:fields-dimension}
Every certified quartic with minimum zero in one or two variables is a sum of
squares of rational polynomials. Consequently, among certified quartics with
minimum zero, the ternary quartics of Theorem~\ref{thm:fields-prime} with
$\ell=5$ and of Example~\ref{ex:fields-ternary} have the smallest possible
number of variables for a failure of rational descent.
\end{corollary}

\begin{proof}
A full positive definite Hessian Gram has a positive definite tensor block
$C$, and comparing the terms of degree two in $X$ gives
$12f_4(X)=(X\otimes X)^{\mathsf T}C(X\otimes X)>0$ for $X\ne0$, where $f_4$
is the quartic part; so the quartic form is positive definite and the degree
is four. The minimizer $p$ is the only real zero, by strong convexity. In one
variable the minimal polynomial of $p$ divides both $f$ and $f'$, hence its
square divides $f$, so $p$ has degree at most two; a real quadratic
irrational zero would bring its real conjugate as a second zero. Thus $p$ is
rational, and Lemma~\ref{lem:taylor-sos}(e) gives a sum of squares over
$\Q$. In two variables $f$ is nonnegative, has exactly one real zero, and
has a positive definite quartic form, so
Lemma~\ref{lem:algebraic-bivariate-descent} applies.
\end{proof}

Lemma~\ref{lem:algebraic-bivariate-descent} rests on Scheiderer's
classification of rational nonnegative ternary quartic forms that are not
sums of squares of rational forms \cite[Theorem~4.1]{Scheiderer2016}; such a
form has at least two real projective zeros, whereas the homogenization of a
certified bivariate quartic with minimum zero has exactly one. The corollary
concerns the strict Hessian hypothesis; we make no minimality claim for
weaker convexity hypotheses.

The obstruction in Theorem~\ref{thm:fields-prime} and
Example~\ref{ex:fields-ternary} is one of positivity: both quartics lie in
the span $\mathcal W^\Q(p)$ of products of rational vanishing quadratics,
but their unique or restricted Gram matrices are indefinite. A cruder
obstruction is also possible: a certified quartic can lie outside that span
altogether. This is the content of the next proposition, whose proof is in
Appendix~\ref{app:fields-four}.

\begin{proposition}[A linear obstruction in four variables]
\label{prop:fields-four}
Let $\alpha=2^{1/3}$, $\beta=5^{1/3}$, and $p=(\alpha,\alpha^2,\beta,\beta^2)$
in the variables $(x,y,z,w)$. There is a rational quartic $f$ with a rational
positive definite full Hessian Gram, $\nabla^2f\succeq I$, and minimum zero at
$p$, such that $f\notin\mathcal W^\Q(p)$. In particular $f$ is not a sum of
squares of rational polynomials. It is a sum of squares over
$\Q(\alpha,\beta)$, a field of degree nine.
\end{proposition}

Here $\dim_\Q\mathcal I_2^\Q(p)=6$, and no product of two rational vanishing
quadratics contains the monomial $yz^3$, whereas the stationary quartic
$y(z^3+w^3/5-3zw+5)$ does. The proof adds a large multiple of a strictly
convex rational sum of squares to this quartic. The general principle is
that, whenever a sum of squares of rational quadratics vanishing at $p$ has a
positive definite full Hessian Gram, any rational stationary quartic outside
$\mathcal W^\Q(p)$ can be made strictly convex in this way, so the quotient
$\mathcal J_4^\Q(p)/\mathcal W^\Q(p)$ is itself an obstruction to rational
descent. Theorem~\ref{thm:fields-descent} below shows that, when a convex
rational sum-of-squares baseline exists, these two obstructions, a large
quadratic vanishing space and a stationary quartic outside the product span,
are the only possible sources of a counterexample.

\subsection{An exponentially large least field}
\label{sec:fields-tower}

The prime family needs dimension proportional to the field degree. A tower of
fifth roots gives exponential degree in polynomial size. Fix $k\ge1$, put
$n=3k$, and let
\[
 a_1=2^{1/5},\qquad a_i=a_{i-1}^{1/5}=2^{1/5^i}\quad(2\le i\le k),\qquad
 p_k=(a_i,a_i^2,a_i^3)_{i=1}^k .
\]
Use variables $(x_i,y_i,z_i)$ for gate $i$, put $b_1=2$ and $b_i=x_{i-1}$
for $i\ge2$, and define the five rational quadratics of gate $i$:
\begin{equation}
 r_{i,1}=x_i^2-y_i,\quad r_{i,2}=x_iy_i-z_i,\quad
 r_{i,3}=y_i^2-x_iz_i,\quad r_{i,4}=y_iz_i-b_i,\quad
 r_{i,5}=z_i^2-b_ix_i .
 \label{eq:fields-tower-relations}
\end{equation}
All of them vanish at $p_k$, since $a_i^5=b_i(p_k)$. The three \emph{chain
residuals} $r_{i,1},r_{i,2},r_{i,4}$ of all gates have $p_k$ as their only
common real zero: they force $y_i=x_i^2$, $z_i=x_i^3$, and $x_i^5=b_i$, and
an odd root of a real number is unique. The relation $r_{k,3}$ of the last
gate is the one that will be subtracted.

\begin{theorem}[An exponentially large least field]
\label{thm:fields-tower}
For every $k\ge1$ one can construct, in time polynomial in $k$, a rational
quartic $F_{k,0}$ that is an explicit sum of $3k+1$ squares of rational
quadratics vanishing at $p_k$, a rational positive definite full Hessian Gram
$M_k$ of $F_{k,0}$, and a positive integer $\lambda_k$, all of total binary
length polynomial in $k$, with the following properties for every rational
$\lambda\ge\lambda_k$ and
\[
 f_{k,\lambda}=\lambda F_{k,0}-r_{k,3}^2 .
\]
\begin{enumerate}[label=(\alph*)]
\item $f_{k,\lambda}$ has a rational full Hessian Gram at least $I$; hence
  $\nabla^2f_{k,\lambda}\succeq I$, the unique minimizer is $p_k$, and
  $\min f_{k,\lambda}=0$.
\item For every subfield $E\subseteq\R$ the following are equivalent:
  $2^{1/5^k}\in E$; $f_{k,\lambda}$ is a sum of squares over $E$;
  $f_{k,\lambda}$ has a positive semidefinite Gram matrix over $E$.
\end{enumerate}
In particular the least real field of both certificate types is
$\Q(2^{1/5^k})$, of degree $5^k=5^{n/3}$ over $\Q$.
\end{theorem}

The baseline has the form
$F_{k,0}=\sigma_0G_k^2+\sigma_1\sum_{i=1}^k\sum_{j\in\{1,2,4\}}r_{i,j}^2$
with positive rationals $\sigma_0,\sigma_1$ and a rational quadratic $G_k$
vanishing at $p_k$, whose homogeneous quadratic part is positive definite and
whose coefficient of $z_k^2$ is positive. Appendix~\ref{app:fields-tower}
constructs it from an exact exposing identity for one quintic gate, a lower
bound for the Jacobian of the chain residuals, and the generic assembly of
Lemma~\ref{lem:quartic-realization}. We sketch the field argument, which is
the heart of the theorem.

Suppose $a_k\notin E$ and put $L=E(a_k^5)$. By
Lemma~\ref{lem:fields-tools}(d), $T^5-a_k^5$ is irreducible over $L$, and all
earlier coordinates $a_i=a_k^{5^{k-i}}$, $i<k$, lie in $L$. Substitute the
earlier coordinates of $p_k$ into a putative certificate over $E$. The result
is a certificate over $L$ for the three-variable slice $f^\flat$ of
$f_{k,\lambda}$ through the last gate. Its five relations
$r^\flat=(r^\flat_{k,1},\ldots,r^\flat_{k,5})$ form a basis of the
$L$-quadratics vanishing at $(a_k,a_k^2,a_k^3)$, and their fifteen products
are linearly independent, because their leading forms
$x^2,xy,y^2-xz,yz,z^2$ multiply to a basis of all ternary quartic forms. All
earlier residuals vanish on the slice, so
\[
 f^\flat=\lambda\Bigl(\sigma_0g^2+\sigma_1\sum_{j\in\{1,2,4\}}
   (r^\flat_{k,j})^2\Bigr)-(r^\flat_{k,3})^2,
\]
where $g=\sum_jc_jr^\flat_{k,j}$ is the restriction of $G_k$. Product
independence makes this Gram matrix on $r^\flat$ unique. The vector
$(0,0,c_5,0,-c_3)$ is orthogonal to the coefficient vectors of $g$ and of $r^\flat_{k,1},r^\flat_{k,2},r^\flat_{k,4}$, and the Gram value
on it is $-c_5^2$. Since only $r_{k,5}$ contains $z_k^2$, $c_5$ is the
positive $z_k^2$ coefficient of $G_k$. A sum of squares or a positive
semidefinite Gram over $L$ would give a positive semidefinite Gram on
$r^\flat$, by Lemma~\ref{lem:fields-tools}(b). This contradiction proves
necessity; Lemma~\ref{lem:taylor-sos}(e) proves sufficiency.

A joint field containing all coefficients does not by itself force any one
coefficient to be complicated; a certificate could in principle spread the
field over many coefficients of small degree. The tower excludes this.

\begin{theorem}[One coefficient of exponential degree]
\label{thm:fields-individual}
Let $\ell$ be an odd prime, $k\ge1$, and $D=\ell^k$. If finitely many
algebraic numbers $\beta_1,\ldots,\beta_t$ generate a field containing
$2^{1/D}$, then $[\Q(\beta_i):\Q]\ge D$ for some $i$. Consequently, in every
sum of squares of $f_{k,\lambda}$ with algebraic coefficients, and in every
positive semidefinite Gram matrix of $f_{k,\lambda}$ with algebraic entries,
some single coefficient or entry has degree at least $5^k$ over $\Q$. Writing
its minimal polynomial densely requires at least $5^k+1$ rational
coefficients.
\end{theorem}

The proof in Appendix~\ref{app:fields-individual} combines the real-field
lemma with a Galois argument. The splitting field of $T^D-2$ has an
automorphism of order $D$ fixing the cyclotomic field, while every element of
a product of symmetric groups of degrees less than $D$ has order with
$\ell$-adic valuation below $k$. The input $f_{k,\lambda}$ and its Hessian
certificate have length polynomial in $k$ and the binary length of
$\lambda$. For $\lambda=\lambda_k$, or any $\lambda$ of binary length
polynomial in $k$, the forced dense output is therefore superpolynomial in
the input length and exponential in the number of variables. It is not
claimed to be exponential in the full input bit length.

The obstruction concerns dense algebraic output of polynomial certificates.
It does not extend to compact representations.

\begin{theorem}[Compact exact certificates for the tower]
\label{thm:fields-compact}
Let $k\ge1$, let $\lambda\ge\lambda_k$ be rational, and let $L_\lambda$
denote $k$ plus the binary length of $\lambda$; for $\lambda=\lambda_k$ this is
polynomial in $k$. Use the signed decomposition
$f_{k,\lambda}=\sum_{j=1}^{3k+2}\sigma_jg_j^2$ with rational weights
$\sigma_j$ and the rational quadratics $g_j$ from
Theorem~\ref{thm:fields-tower}: the exposing quadratic $G_k$, the $3k$ chain
residuals, and $r_{k,3}$; the coefficient of $r_{k,3}^2$ is $-1$.
\begin{enumerate}[label=(\alph*)]
\item \emph{Shared root circuit.} In time polynomial in $L_\lambda$ one can
  construct an arithmetic circuit of size polynomial in $L_\lambda$ whose only
  nonrational
  operations are $k$ real fifth-root gates computing $a_1,\ldots,a_k$, and
  polynomially many quadratics whose coefficients are outputs of the circuit,
  such that $f_{k,\lambda}$ is the sum of their squares.
\item \emph{Auxiliary root variables.} In time polynomial in $L_\lambda$ one
  can construct a rational sum of squares $s(X,Y)$ of degree four in $6k$
  variables and rational polynomials $H_j(X,Y)$ of degree at most two, of
  total size polynomial in $L_\lambda$, with
  \[
   f_{k,\lambda}(X)=s(X,Y)+\sum_{j=1}^{3k+2}H_j(X,Y)\,g_j(Y).
  \]
  The equations $g_j(Y)=0$ have the real solution $Y=p_k$, so the identity
  proves $f_{k,\lambda}\ge0$. Eliminating the two redundant equations
  $G_k(Y)=0$ and $r_{k,3}(Y)=0$ gives an identity of degree at most five
  modulo the $3k$ chain residuals alone.
\end{enumerate}
\end{theorem}

Part (a) is Taylor integration with symbolic coordinates
(Lemma~\ref{lem:taylor-sos}(e)); part (b) is the first-order convexity
identity of SOS-convex polynomials \cite[Theorem~3.1]{AhmadiParrilo2013}
written with the minimizer as a variable. Neither is a new proof system.
The point of the theorem is the comparison on one family: the input and its
convexity certificate are short; dense algebraic polynomial certificates of
the optimal value are exponentially long; shared root circuits and rational
certificates modulo the root equations are short. The existence of a real
solution in (b) is essential, since an identity modulo an inconsistent
system proves nothing. Rational-function certificates give another short
alternative (Corollary~\ref{cor:fields-tower-denominator}). The theorem gives
no lower bound for sparse polynomials, radicals, root towers, or arithmetic
circuits; $T^{5^k}-2$ itself is sparse.

\subsection{Rational relations and stationary quartics at the tower point}
\label{sec:fields-spaces}

The least-field proof uses one direction of one slice. We now describe all
rational quadratic relations and all rational stationary quartics at the
tower point; the recognition corollary then applies the required positive
semidefinite test.

\begin{theorem}[Rational relations and stationary quartics]
\label{thm:fields-tower-spaces}
Let $k\ge1$ and let $r$ be the vector of the $5k$ quadratics
\eqref{eq:fields-tower-relations}.
\begin{enumerate}[label=(\alph*)]
\item The entries of $r$ form a basis of $\mathcal I_2^\Q(p_k)$, and no
  nonzero rational affine polynomial vanishes at $p_k$. Thus
  $\dim_\Q\mathcal I_2^\Q(p_k)=5k$.
\item The $\binom{5k+1}2$ products $r_{i,j}r_{i',j'}$ are linearly
  independent over $\R$.
\item $\mathcal J_3^\Q(p_k)=\{0\}$ and
  $\mathcal J_4^\Q(p_k)=\mathcal W^\Q(p_k)
   =\{r^{\mathsf T}Sr:S\in\operatorname{Sym}_{5k}(\Q)\}$, with $S$ unique.
\end{enumerate}
\end{theorem}

The proof is in Appendix~\ref{app:fields-spaces}. Part (a) eliminates one
pivot monomial per relation and identifies the remaining monomials with
distinct base-five digit patterns of exponents of $a_k$ below $5^k$. Part (b)
inducts over the gates, starting from the fifteen leading products of one
gate. Part (c) uses a first-order interpolation lemma for one quintic gate,
whose matrix splits into five blocks according to the weight of a monomial
modulo five, and an induction over gates that removes the last gate's
contribution. All arguments are uniform in $k$.

\begin{corollary}[Exact rational recognition at the tower point]
\label{cor:fields-recognition}
Given $k$ and a rational polynomial $f$ of degree at most four, one can
decide in time polynomial in $k$ and the input length whether $f(p_k)=0$
and, if so, whether $f$ is a sum of squares of rational polynomials. The
rational positive semidefinite Gram matrix of such an $f$ on $z$ is unique,
and in the affirmative case a rational sum of squares of polynomial size is
produced. If $f(p_k)=0$ and $\nabla f(p_k)=0$, for instance when $p_k$ is an
unconstrained minimizer with value zero, the linear step that writes
$f=r^{\mathsf T}Sr$ always succeeds, and $f$ is then a sum of squares of
rational polynomials exactly when $S\succeq0$; the tower quartics
$f_{k,\lambda}$ are stationary at $p_k$ and have $S\not\succeq0$.
\end{corollary}

\begin{proof}
Evaluate $f(p_k)$ by writing each monomial value as $2^ta_k^e$ with
$0\le e<5^k$; since $1,a_k,\ldots,a_k^{5^k-1}$ are linearly independent over
$\Q$, $f(p_k)=0$ exactly when the coefficients collected in each residue
class sum to zero. Exponents have $O(k)$ bits, and no dense field
representation is formed. If $f$ is a rational sum of squares, its factors
lie in $\mathcal I_2^\Q(p_k)$ by Lemma~\ref{lem:fields-tools}(a,b), so
$f=r^{\mathsf T}Sr$ with $S\succeq0$ rational. Conversely such a
representation gives rational squares (Lemma~\ref{lem:models-rational-squares}).
By Theorem~\ref{thm:fields-tower-spaces}(b) the matrix $S$ is unique and is
found by solving a rational linear system with $O(k^2)$ unknowns and $O(k^4)$
equations whose matrix has bounded integer entries; one rational symmetric
elimination tests $S\succeq0$. For a rational positive semidefinite $Q$ with
$f=z^{\mathsf T}Qz$, Lemma~\ref{lem:fields-tools}(b) gives
$Q=BSB^{\mathsf T}$ with $r=B^{\mathsf T}z$, so $Q$ is unique as well. The
last sentence follows from Theorem~\ref{thm:fields-tower-spaces}(c), which
places every rational quartic with zero value and gradient at $p_k$ in the
span of the products of relations, and from the slice argument for
Theorem~\ref{thm:fields-tower}(b), which shows that the unique matrix of
$f_{k,\lambda}$ is indefinite.
\end{proof}

The recognition is exact and avoids both a semidefinite feasibility problem
with a free matrix and any search over algebraic coefficient fields. Its
scope is narrow: it requires the zero to be the supplied tower point, it
decides rational, not real, sums of squares, and it does not locate an
unknown zero. The theorem also explains the tower construction. Every
rational stationary quartic at $p_k$ is a quadratic form in the relations,
so a strictly convex real sum of squares can fail rational descent only
through an indefinite unique matrix $S$, which is exactly what
Theorem~\ref{thm:fields-tower} produces. The rational span of the relations
is much smaller than the real space $\mathcal I_2^\R(p_k)$, which has
dimension $\binom{3k+2}2-1$; real square factors may use all of it.

\subsection{A descent criterion}
\label{sec:fields-criterion}

The constructions above subtract the square of a rational vanishing quadratic
that the convex baseline does not use. This requires more rational vanishing
quadratics than the dimension forces. The following criterion shows that
this construction is impossible when the quadratic vanishing space has its
minimum possible dimension. It uses the real sum-of-squares length bound of
Theorem~\ref{thm:sos-length}: a globally convex real quartic of degree
exactly four with a zero at which its Hessian is positive definite is not a
sum of fewer than $n+1$ squares of real polynomials.

\begin{theorem}[Descent from a small vanishing space]
\label{thm:fields-descent}
Let $p\in\R^n$ satisfy
\[
 \dim_\Q\mathcal I_2^\Q(p)=n+1
 \quad\text{and}\quad
 \mathcal J_4^\Q(p)=\mathcal W^\Q(p),
\]
and suppose there is a globally convex rational polynomial $G$ of degree
exactly four that is a sum of squares of rational polynomials and satisfies
$G(p)=0$ and $\nabla^2G(p)\succ0$. Let $g$ be a rational basis of
$\mathcal I_2^\Q(p)$. Then every globally convex rational polynomial $F$ of
degree exactly four with
\[
 F(p)=0,\qquad \nabla F(p)=0,\qquad \nabla^2F(p)\succ0
\]
has a unique symmetric matrix $S$ with $F=g^{\mathsf T}Sg$, and $S$ is
rational and positive definite. In particular $F$ is a sum of squares of
rational polynomials.
\end{theorem}

\begin{proof}
By Lemma~\ref{lem:fields-tools}(a,b), the rational square factors of $G$ lie
in $\mathcal I_2^\Q(p)$, so $G=g^{\mathsf T}S_0g$ with $S_0\succeq0$ rational.
If $S_0$ were singular, a real factorization of $S_0$ would write $G$ as a sum
of at most $n$ real squares, contradicting Theorem~\ref{thm:sos-length}. Thus
$S_0\succ0$. Since $F\in\mathcal J_4^\Q(p)=\mathcal W^\Q(p)$, there is a
rational symmetric $S$ with $F=g^{\mathsf T}Sg$. Suppose $S$ is not positive
definite. Along $S_t=(1-t)S_0+tS$ the least eigenvalue is continuous,
positive at $t=0$ and nonpositive at $t=1$, so there is a first
$t_*\in(0,1]$ at which $S_{t_*}$ is positive semidefinite and singular. Then
$F_{t_*}=(1-t_*)G+t_*F=g^{\mathsf T}S_{t_*}g$ is a sum of at most $n$ real
squares. On the other hand $F_{t_*}$ is globally convex, vanishes at $p$, and
has positive definite Hessian there. It also has degree exactly four. The
quartic parts $G_4$ and $F_4$ are nonnegative, being limits of
$t^{-4}G(tx)$ and $t^{-4}F(tx)$ and $F\ge F(p)=0$ by convexity and
stationarity; if $t_*<1$ and $(1-t_*)G_4+t_*F_4=0$ then $G_4=0$, which is
false, and if $t_*=1$ the degree of $F$ is assumed. This contradicts
Theorem~\ref{thm:sos-length}, so $S\succ0$, and
Lemma~\ref{lem:models-rational-squares} gives rational squares.

For uniqueness, let $K\ne0$ be a real symmetric matrix with
$g^{\mathsf T}Kg=0$; replacing $K$ by $-K$ if necessary, $K$ has a negative
eigenvalue. The ray $S_0+tK$, $t\ge0$, reaches a first singular positive
semidefinite matrix at a finite $t>0$, which represents $G$ by at most $n$
real squares. This is again impossible.
\end{proof}

The intermediate parameter $t_*$ need not be rational; this is harmless,
because Theorem~\ref{thm:sos-length} allows real square factors. The proof
does not assume that the individual factors are convex.

\begin{corollary}[Product independence]
\label{cor:fields-product-independence}
If $\dim_\Q\mathcal I_2^\Q(p)=n+1$ and a baseline $G$ as in
Theorem~\ref{thm:fields-descent} exists, then the products of a basis of
$\mathcal I_2^\Q(p)$ are linearly independent over $\R$ and
$\dim_\Q\mathcal W^\Q(p)=(n+1)(n+2)/2$. Under all hypotheses of
Theorem~\ref{thm:fields-descent}, the quartic $F$ has a unique rational
positive semidefinite Gram matrix on $z$, of rank $n+1$, and is a sum of
exactly $n+1$ real squares and of no fewer.
\end{corollary}

\begin{proof}
The uniqueness part of the proof above used only $G$ and the dimension. For
the Gram statement, Lemma~\ref{lem:fields-tools}(b) gives
$Q=BS'B^{\mathsf T}$ for any rational positive semidefinite Gram $Q$ of $F$,
where $g=B^{\mathsf T}z$, and uniqueness gives $S'=S$. A real factorization of
$S\succ0$ gives $n+1$ squares; Theorem~\ref{thm:sos-length} gives the lower
bound.
\end{proof}

The two hypotheses have different roles. The dimension bound turns a
singular restricted Gram into too few real squares; the equality
$\mathcal J_4^\Q(p)=\mathcal W^\Q(p)$ supplies a restricted Gram for every
rational stationary quartic. Example~\ref{ex:fields-ternary} has
$\dim\mathcal I_2^\Q=5>n+1=4$, Proposition~\ref{prop:fields-four} violates
the second hypothesis, and the tower has $\dim\mathcal I_2^\Q(p_k)=5k>3k+1$.
A counterexample to rational descent at a point carrying a rational convex
sum-of-squares baseline must therefore use more than $n+1$ rational
vanishing quadratics or a stationary quartic outside their product span. The
criterion is an obstruction to a construction method; it gives no lower
bound for other certificate systems.

The dimension hypothesis holds for the cyclic quartics of
Theorem~\ref{thm:algebraic-cyclic}, whose zeros have field degree
exponential in $n$.

\begin{proposition}[Vanishing quadratics at the cyclic points]
\label{prop:fields-cyclic}
Let $n\ge4$, and let $m=n+1$, $\eta=(-1)^m$, $d_n=(2^m-\eta)/3$,
$e_i=((-2)^i-1)/3$, and $p=(p_1,\ldots,p_n)$ with $p_i=2^{e_i/d_n}$, as in
Theorem~\ref{thm:algebraic-cyclic}. With $x_0=1$ and indices modulo $m$, the
$m$ cyclic binomials $x_i^2-2^{\theta_i}x_{i+1}x_{i+2}$, where $\theta_i=0$
except $\theta_{m-2}=\eta$ and $\theta_{m-1}=-\eta$, form a basis of
$\mathcal I_2^\Q(p)$. Thus $\dim_\Q\mathcal I_2^\Q(p)=n+1$, and by
Corollary~\ref{cor:fields-product-independence} applied to the rational
sum-of-squares quartic of Theorem~\ref{thm:algebraic-cyclic}, the products of
these relations are linearly independent.
\end{proposition}

The proof is in Appendix~\ref{app:fields-cyclic}. For $n=3$ the space is
larger, five-dimensional, and the hypothesis fails. Whether
$\mathcal J_4^\Q(p)=\mathcal W^\Q(p)$ holds at all cyclic points is open.
If the equality holds for a given $n\ge4$, then
Theorem~\ref{thm:fields-descent} shows that no globally convex rational
quartic with minimum zero at that cyclic point and a positive definite
Hessian there can fail rational
descent, although the point has algebraic degree $d_n$.
